\scanpage{42}\mspage{19}

$u\in M$ et pour $p,q\in\mathbb N$, $\varphi^pu=\varphi^qu\Longrightarrow p=q$,
ou encore \deleted{que $\varphi^p$} (se ramenant au cas $u=\varphi$
puisque $\deleted{\unsure{\exists}}u'$ avec $uu'=\varphi$) $\varphi^p=\varphi^q\Longrightarrow p=q$.
Dans ce cas, par des principes généraux,
\deleted{le quotient} l'objet $e$ de $X^\wedge$ étant $\unsure{0|2}$-connexe,
l'objet quotient $F=e/\mathbb N=e/\varphi$ est $0$-connexe,
et son $\pi_1$ admet un groupe quotient isomorphe
à $\mathbb Z$, donc $\deleted{\unsure{F}}X_{/F}$ n'est pas $1$-connexe,
et $=\ H^1(F,\deleted{\unsure{Q|\mathbb Z}}k)\ne0$ pour tt groupe
abélien $k\ne0$.

b) Il existe $p\ne q$ tel que $\varphi^p=\varphi^q$. Alors
\deleted{l'opération} l'image de $M$ dans l'ens.
$\unsure{\mathrm{Ep}|\mathrm{E}p}(E)$ des épimorphismes de $E=$ ens des
éléments minimaux de $E$ est un \uline{groupe}
$\Gamma$, et si $X_0$ désigne l'ouvert de $X$
formé des plus petits objets de $X$, alors
\[
\deleted{\unsure{\pi}}\quad \pi_1(X_0,e)\simeq\pi_1(M)\xrightarrow{\text{épi}}\Gamma.
\]
Donc si $X_0$ est $1$-connexe, alors $\Gamma=\{1\}$, donc
on a $u\varphi=\varphi\quad\forall u\in M$, et $M$, \deleted{et} donc
aussi $X_{\deleted{\unsure{0}}}$, \deleted{\unsure{et $X$}} sont cofiltrants. \deleted{\illegible{1}}
% slanted
\margin{En particulier, on a $\varphi^2=\varphi$ pour tt $\varphi$ minimal.}

\uline{Corollaire} Si \unsure{tout|$\mapsto$|tous les} objet de $X^\wedge$ est $1$-connexe,
alors $X$ est cofiltrante,

% archive sheet number: 41
